\chapter*{如何阅读本手册}
\addcontentsline{toc}{chapter}{如何阅读本手册}
InkStruct 编辑器将 Markdown 源码、实际手写轨迹和毫米页面结合起来。左侧写内容，右侧检查排版；正式导出前，程序检查缺字、尺寸与可书写内容。用系统字体显示在输入框中的文字，不等于已经拥有它的手写字形。

本手册依据当前编辑器界面和实现编写。图示采用本项目的实际测试界面；示例数字仅用于学习操作，不代表实验测量结果。手册中的普通印刷字体用于讲解，与应用导出的手写资源不同。

\section*{按任务选择阅读路径}
\par\noindent\begin{tabularx}{\linewidth}{@{}P{36mm}Y@{}}
\toprule
我要完成的任务 & 建议阅读\\\midrule
第一次打开应用 & 第\ref{ch:basics}章认识文件，第\ref{ch:quick}章完成数字示例。\\
排版一份报告 & 第\ref{ch:workspace}至\ref{ch:math}章，按需要加入表格与公式。\\
在预印表单上填写 & 第\ref{ch:templates}章，先建模板，再绑定内容。\\
检查生成字形好坏 & 第\ref{ch:resources}章，导入资源并使用字形评价。\\
交给写字机 & 第\ref{ch:export}章，区分参考页面与机器文件。\\
让 Agent 帮忙设计模板 & 附录\ref{ch:agent}，交付模板、项目与适配报告。\\
\bottomrule
\end{tabularx}

\section*{统一手册的组织方式}
全书按四个分部组织。第一部放置所有模块共享的项目与资源概念；第二部完整介绍当前编辑排版；第三、四部分别介绍生成和采集的操作与交接。后续扩展在各分部中增加操作章节，不将生成参数、设备校准等内容塞入编辑器章节。

\note{统一应用通过左侧导航提供采集、生成和排版，底部依次为样本库、字体库、模板库和设置。生成默认 v6.2 exp，v6.1 exp 和 v4 adaptive 可选；具体步骤见第\ref{ch:generation}章。}

\section*{约定}
\ui{方括号式文字}表示界面按钮、菜单或选项；等宽框展示可输入的源码。尺寸默认使用毫米。macOS 快捷键以 Command 表示，Windows 以 Ctrl 表示；有差异时以应用菜单显示为准。窗口图示来自 macOS/Qt，Windows 的原生对话框外观可能不同。

本册源码适用应用 0.1.88，部分历史图示沿用 macOS 外观。Windows 启动时使用独立工作区，在界面重新选择本机样本目录。Windows 的实际验收结果以项目内的 Windows 初始化记录为准。真实数位板、驱动、中文输入法组合输入和写字机纸面效果仍需实物验证。

\paragraph{历史排版更新（0.1.43）}本次主应用合入两端对齐、中文自然书写v2与排版问题面板。自然书写默认关闭，旧v1效果保留；显式强制SVG只输出已有有效笔迹并附缺口报告，不代表文稿完整通过常规导出检查。应用正文参数时保留本文模板和纸张，明确修改对应项才替换。历史独立0.1.43预览ZIP不包含打包后新增的自然书写v2与问题面板，应以本次主应用包及交付记录为准。

\paragraph{大字体与成品数量（0.1.46）}字体展开总量上限为4 GiB，按成员未压缩大小计算，不是压缩文件体积；单成员、点数和完整校验仍保留。v6.2生成页的\ui{每字成品上限}可选1--64份，默认4份，包含原主字形；真实供体按完整部件均衡尝试，同一供体字可使用不同录制样本，也可有限尝试两个部件同时变化，逐份检查最终整字。来源不足时不凑数。旧字体、项目和冻结选样不自动改写，方法名称不变。详见第\ref{ch:generation}章与第\ref{ch:resources}章。

\paragraph{生成进度（0.1.57）}浅绿区域的宽度表示已处理比例，深绿区域显示各阶段的处理速度。当前速度以字/秒显示，读数位于图内右侧并随横线升降。图下可查看剩余字数、失败数及预计剩余时间。

\paragraph{长文稿预览（0.1.12）}字体候选和字形边界可在内容完全一致时复用；自动恢复副本在后台保存，期间可继续输入。正式保存仍检查文件是否被其他程序修改。缺字和未支持公式仍会阻止成品导出，性能优化不改变此限制。

\paragraph{连续页面预览（0.1.13）}右侧页面纵向连续排列，可直接滚轮浏览；页码随当前位置更新，输入页码可跳转。缩放时保持当前阅读位置，点击其他页的字形仍可定位源码。屏幕图像按显示缩放缓存，矢量导出继续使用原始笔迹。首次进入未缓存页面时会短暂显示绘制提示；快速滚动仅优先绘制当前位置，内存缓存有上限。

\paragraph{统一设置（现行位置0.1.67）}左下角设置是采集与排版设置的统一入口，支持 Ctrl+逗号。应用设置的主题提供跟随系统、浅色、深色，保存后立即生效。公式设置的两个对照预览采用当前主题配色。采集显示只影响屏幕。汉字与标点扩展的笔迹粗细在采集页分别调整。排版的 Ctrl+N 从全局默认新建文档，导入 Markdown 时自动嵌入源码目录内受支持的本地图片；排版设置保存在当前项目。模板和字体默认值分别在模板库和字体库管理。设置左侧按正文、公式、表格、图片、代码和脚注分类；修改先保留为草稿；存为默认只影响新文档并关闭窗口，应用并保存同时更新当前文档且保持窗口打开，取消放弃尚未保存的修改。表格和图片可沿用文档默认或独立设置，修改默认不覆盖独立设置。问题诊断与字形评价仍在排版工作区。

\paragraph{固定浅色界面（0.1.15）}当前版本固定使用浅色界面，不跟随系统深色模式。此设置仅作用于应用显示，不更改操作系统主题，也不改变文档、笔迹或字体参数。深色模式兼容留待后续版本。

\paragraph{全局默认与元素属性（0.1.16）}齿轮保存本机全局默认，供后续新文档使用；已有项目保留原参数。可明确勾选同时用于当前文档。排版右侧的“元素属性”用于选择表格或图片、设置独立参数或恢复沿用本文默认，点击应用后随项目保存。修改全局默认不清除已有独立属性。公式设置窗口按内容最小宽度限制，支持纵向滚动。

\paragraph{字体导入进度（0.1.17）}导入时先复制到临时区域，完整验证后加入字体库，显示当前阶段与已验证候选数。完整字体仍需校验每个成员、轨迹、动态信息、来源和度量；已有同包编号但内容不同的文件会拒绝导入并保留原文件。

\paragraph{并行验证与防误调（0.1.18）}完整字体导入最多使用四个进程分批校验，小字体和低核心环境串行处理，完整检查不变。参数框、参数滑块和下拉框不使用滚轮改值，经过控件时仍可滚动页面；点击和键盘调整保留。正文设置全部展开，公式页移除旧容器箭头，实际字形预览默认显示辅助线。以上界面行为不改写旧文档参数。

\paragraph{字符校正（0.1.20）}表格、图片设置直接陈列全部选项。点击排版预览字符，在右侧调整位移和宽高比例。范围为“仅此处”“本文同一样本”“本文同一字符的全部样本”：例如选中 A2，依次只改此处、全部 A2 或同类 A1/A2 等所有样本。正斜体默认分开，跨类别需明确勾选；界面显示样本及影响次数。批量操作统一更新范围内的校正，旧单字校正不会遮住结果，可整体撤销、重做或恢复。后两种范围可勾选“同时保存到字体包的新版本”，完整校验并保留原包和来源；成功后应用本文校正，当前文档使用原字体快照，不重复叠加，新包可导入后用于后续文档。取消或失败不修改本文，处理期间文档变化不会套用旧操作。校正随项目保存并用于 SVG，字距、换行、原资源和超界检查保留；旧文档按原效果重开。

\paragraph{模板连续正文（0.1.21）}简单模板重复单页正文，自定义模板使用有序起始页组与循环页组。一个 Markdown 同时存放命名字段与唯一正文，正文按编号跨框流动，字段独立固定区。模板制作仍从齿轮进入，本份文档从右侧“模板内容”套用和填写。旧模板及项目保持原版读取。

\paragraph{模板界面整理（0.1.22）}齿轮中的模板管理与排版默认分开。工作台统一草稿编辑，纸张、框与字段使用明确属性；右侧“模板”提供字段表单和适配问题。表单仍修改唯一 Markdown，保留原协议和冻结选样。

\paragraph{英文字距与邻字参考（0.1.23）}正文与段落设置增加按笔迹宽度排列英文字母，排除单字采集的横向位置留白，不缩放笔迹。旧文档和已有默认保持旧字距；新安装使用推荐值。明确勾选将默认排版参数应用到当前文档时，会同时应用设置中的其他排版默认，应先核对。字符校正预览显示同一行真实邻字及字距，淡蓝底标识选中位置，灰色为原字形，蓝色为校正草稿；批量范围同步预示可见邻字变化，应用前不改变文档与冻结选样。

\paragraph{图片高度与汉字公式（0.1.24）}图片默认高 60 mm、宽度按比例，默认页不再提供宽度输入。单图可独立选择尺寸依据。旧文档宽度与个性化属性继续保留，转换为新默认需明确操作。公式直接支持汉字及汉字上下标，例如 \code{U_{模块}}，仍检查真实字体、源码与导出能力。

\paragraph{跳过参考图片导出（0.1.25）}含非 SVG 图片时先提示再导出；确认后保留占位与图题，输出清单记录跳过项。图片资源与文档保留，缺字、公式及越界检查继续生效。

\paragraph{缺失位图引用的导出（0.1.26）}尚未载入的本地 PNG/JPEG 可与已载入位图一起明确确认跳过。弹窗分别计数并可查看文件名；缺失错误保留，其他阻断问题仍需修正。导出保持当前占位、图题、页序、坐标和冻结选样。

\paragraph{模板预览停止（0.1.28）}某对象无法放入整个循环页组时，保留此前已排正文及已引用脚注，明确提示此处及后续正文尚未排版。通过查看问题定位原稿，调整正文框或手动拆分公式；不自动缩放、截断或允许成品导出。较大字体资源目录采用独立有界容量，校验与原文档保留。

\paragraph{标点间距与路径优化（0.1.29）}新文档的中英文标点均按实际笔迹宽度紧凑排列，不保留半字或整字空位。保留最小笔迹间隙、大小、基线与选样；公式及代码沿用原规则。齿轮中的正文与段落提供标点间距选择，应用到旧文档须明确勾选。路径优化默认开启，旧文档保留原值；优化只调整允许的轨迹顺序与方向，不改变字形和坐标。

\paragraph{紧凑公式（0.1.30）}新文档公式消除采集位置带来的横向留白，按真实笔迹边界布置分式、根式和括号，降低默认运算符及结构净距，保留笔宽保护和显式空格。根号和伸缩括号默认使用录制字形；根号钩部与上横线分段适配，括号按样本比例和类型限制宽度。旧文档保留原样式，明确应用新默认后才重排；原字体和选样不改写。

\paragraph{预览与 SVG 导出（0.1.33）}左下角设置的“预览绘制”可选择 CPU（默认）或 GPU，保存后下次启动生效；GPU 失败自动回退 CPU。两者均复用有限 SVG 解析缓存，并避免反复取消相邻页预取。只影响显示，不改变排版和选样。
“导出 SVG”打开选项弹窗：页码留空表示全部，也可输入 1-3,5；可取消表格内内容，仅保留表格线和表外内容。路径优化默认开启。先严格校验原始排版，再以左上角为原点同步缩放输出点坐标、SVG 画布宽高与 viewBox，笔宽不变；输出尺寸随比例改变。原文档纸张、预览和冻结选样不变。缩放后的路径及画布仍须通过校验。定位点为各内容文本框左上角、右下角边界上的 0.1 mm 短线，笔宽 0.01 mm；校正比例也作用于短线。保持冻结快照及原文档页码，取消不输出。
制作字体时，来源齐全且目标未录制时可合成摄氏度（中文句号缩小抬高，加正体 C），短横线复用减号；全角波浪号从两笔约等号各提取一笔，每份约等号提供两份候选。记录来源与变换，不伪造独立采集，不覆盖旧字体和冻结文档。

\paragraph{窄数字占位（0.1.32）}紧凑公式对数字保留最小占位，避免接近竖线的数字1与邻字过密。参照表格默认数字参考高度与占位下限，最小占位为当前层公式字号的0.315倍；只对不足部分居中补白，不缩放笔迹。分式和上下标按各自字号计算，非数字和表格普通数字规则保持。旧版采集留白公式不受影响，已有紧凑公式重新排版时直接使用修复规则，原稿、样式和选样保留。

\paragraph{v6.2 exp 验证范围}新增止部关键笔和竖弯钩语义准入，保留 v6.1 与 v4。user\_001 验证合并字体仅重生孔、正，其余字形保持；孔末笔保留明确参考缺口，不代表全作者字体。已有文档继续使用已冻结资源。


\paragraph{录制符号与翻页（0.1.35）}正文星号 U+002A 与数学星号 U+2217 保留独立字符身份，使用减号录制笔迹缩短后按 0、60、120 度组合；优先三份不同样本，不足时明确记录重复使用。确定性循环配对，每个目标至多与合格减号候选数量相同，不枚举所有组合。公式内星号的原解析含义不变。
数学比号 U+2236 缺失时优先复用全角冒号，再使用半角冒号；数学圆圈 U+2218 缺失时复用中文句号，按数学轴放置，圆圈命令的含义和间距类别不变。整份笔迹等比缩放，不描绘理想点或圆。独立目标优先，记录来源、仿射变换、构造身份和时间平移；不是独立录制。旧字体扩展另存新包，不改旧文档与冻结选样。
排版页页码右侧向上按钮跳至上一页，向下按钮跳至下一页；键盘方向一致，可直接输入页码，连续滚动仍同步当前页。

\paragraph{侧栏与预览（0.1.57）}系统标题栏显示 InkStruct 与应用版本，不再占用独立品牌横栏。左下设置位于模板库下方；Ctrl+逗号保留。Alt+1/2/3/4/5切换采集、生成、排版、字体库、模板库；F1或侧栏右键打开关于。排版预览支持 Ctrl＋滚轮缩放。启动时可直接显示空白 PDF 模板的已校验底图，默认字体在后台校验。依赖字体的页码将在校验后绘出。模板工作台可追加或替换目标页组的 PDF 页面版式。字体库预设直接展示汉字或公式笔迹，点击\ui{编辑试写…}输入自定义内容，切换预设后草稿保留。来源、缺字和完整校验状态始终如实显示。

在两台电脑之间录制样本时，可通过「采集 › 书写者资料」导出和导入同一稳定书写者 ID。返回原设备后，将 JSON、ZIP 或样本目录导入已有书写者。旧汉字批次与另一台电脑录制的标点可在制作字体时核对双方 ID 和样本范围，确认后在本机建立批次级关联。

\paragraph{汉字间距（0.1.85）}在「设置 › 正文与段落」可选整字宽或按笔迹紧凑排列。保存项目后保留所选模式，旧项目按整字宽打开。自然书写的字距变化可独立调节。

\paragraph{样本库（0.1.86）}在「采集」打开样本库，按书写者、字符、类别或任务筛选，预览原始笔迹并选择导出。接收端选择目标书写者，核对离线采集包中的新增、重复和冲突后导入。

\paragraph{段落对齐区块（0.1.86）}分别在独立源码行使用 \verb|\begin{paragraphalign}{center}| 和 \verb|\end{paragraphalign}|。对齐方式可选 \texttt{left}、\texttt{center}、\texttt{right} 或 \texttt{justify}，作用于区块中的普通正文。区块可跨页或正文框，错误指令对应的正文沿用文档默认对齐。

\paragraph{字母与数字间距（0.1.87）}在「正文与段落」应用并保存，可让正文数字随字母按实际笔迹宽度排列。旧项目打开时延用原间距。

\paragraph{独立样本库（0.1.88）}从左下样本库或 Alt+6 打开。采集页的快捷入口将带入当前书写者。筛选结果按字符分组，右侧可切换份次并查看笔迹和来源。展开组后可逐份多选和导出。
